\section{Compiler Construction Basics}\label{sec:background-compilers}

A compiler is a specialized software system that translates source code written in a high-level programming
language into a lower-level representation, such as machine code, bytecode, or in the context of this thesis,
a binary MIDI file.
The compilation process is traditionally divided into several distinct phases to manage complexity and ensure
correctness.

\subsection{Lexical and Syntactic Analysis}

The first phase, \textit{Lexical Analysis} (or scanning), reads the raw character stream of the source file
and groups characters into meaningful sequences called \textit{tokens}.
For example, the string \texttt{tempo} might be converted into a \texttt{KEYWORD\_TEMPO} token, while
\texttt{120} becomes an \texttt{INT\_LITERAL} token.
This is often implemented using tools like Flex, which rely on regular expressions.

The second phase, \textit{Syntactic Analysis} (or parsing), takes the stream of tokens and verifies that they
form a valid sequence according to the formal grammar of the language.
This phase organizes the flat token stream into a hierarchical tree structure known as an \textit{Abstract
Syntax Tree} (AST).
The AST captures the logical operations of the code (e.g., nesting a \texttt{play} statement inside a
\texttt{loop} block) while discarding irrelevant textual details like whitespace and semicolons.

\subsection{Semantic Analysis}

While the parser ensures that a sentence is grammatically correct, the \textit{Semantic Analyzer} ensures
that it makes logical sense.
This phase traverses the AST to perform checks that cannot be captured by a context-free grammar.
Key tasks include:
\begin{itemize}
  \item \textbf{Type Checking:} Ensuring that operations are applied to compatible data types (e.g.,
    preventing the addition of a boolean to a note).
  \item \textbf{Symbol Resolution:} Maintaining a \textit{Symbol Table} to ensure that variables and
    functions are declared before they are used, and that they are accessed within their valid scope.
\end{itemize}

\subsection{Intermediate Representation and Code Generation}

After validation, the compiler often translates the AST into an \textit{Intermediate Representation} (IR).
An IR is a lower-level, often flattened representation of the program that is easier to optimize and
translate than the hierarchical AST.
In standard compilers, this might involve generating three-address code.
In the context of music generation, lowering to an IR involves unrolling loops and resolving relative timings
into absolute timestamps.

Finally, the \textit{Code Generation} phase translates the IR into the target output format.
This phase handles the intricate binary serialization details, such as little/big-endian encoding and
specific protocol constraints, yielding the final executable artifact.
